% !TeX root = ../../main.tex
\section{\bilingual{Main Results}{本文主要工作与章节安排}}

    \bilingual{
        The main work of this paper is based on the works of Mallein~\mycite{cortines2016nbranchingrandomwalkrandom}, A\"id\'ekon and Hu~\mycite{aïdékon2025averageweightpercolationcompletegraph}. Using their methods, we study the asymptotic speed of the frontier position of $N$-BRW in the case where the step size has a Gaussian tail with $\alpha=2$. Under some standard assumptions, we prove that the velocity correction term of the frontier position follows a $(\log N)^{-2}$ scaling law with respect to the population size $N$, and give an explicit expression for the constant $C$. This result verifies Brunet--Derrida's conjecture and provides a foundation for understanding the frontier dynamics under finite population selection mechanisms.
    }{
        本文的主要工作,~ 是在 Mallein~\mycite{cortines2016nbranchingrandomwalkrandom},~ A\"id\'ekon 与 Hu~\mycite{aïdékon2025averageweightpercolationcompletegraph} 的基础上,~ 在步长具有$\alpha=2$的Gaussian尾部的情形下, ~使用他们的方法研究$N$-BRW 的前沿位置的渐近速率, ~在一些标准的假设下,~ 证明了前沿位置的速度修正项相对于种群规模 $N$ 的 $(\log N)^{-2}$ 标度律,~ 并且给出了显式的常数 $C$ 的表达式. ~这一结果验证了 Brunet--Derrida 的猜想,~ 为理解有限种群选择机制下的前沿动力学提供了基础.
    }

    \bilingual{
        Before stating the assumptions,~ we introduce some notations that will appear in the sequel: $\Xi$ denotes the point process on $\mathbb R$ that serves as the branching law; $\psi(\lambda):=\log\mathbb E[\sum_{x\in\Xi}e^{-\lambda x}]$ is its log-Laplace transform; $\#\Xi$ is the number of atoms in $\Xi$, and $\Xi_i$ is the $i$-th smallest atom of $\Xi$; $W_n$ is the additive martingale defined via $\psi$ (see Appendix); $\xi$ denotes a random variable whose distribution is given by \nameref{lemma:many-to-one}, referred to as the step of the spine random walk, with variance $\sigma^2:=\psi''(1)$; $\ell_N$ is a time sequence indexed by $N$; $v_N$ is the asymptotic speed of the leftmost particle. Detailed explanations of the assumptions can be found in Section \ref{section:explanations-on-assumptions} of the next chapter:
    }{
        在给出假设前,~ 先说明以下将在后文出现的记号：$\Xi$ 表示作为分支律的 $\mathbb R$ 上的点过程; $\psi(\lambda):=\log\mathbb E[\sum_{x\in\Xi}e^{-\lambda x}]$ 为其对数 Laplace 变换; $\#\Xi$ 为 $\Xi$ 中原子个数, $\Xi_i$ 为 $\Xi$ 的第 $i$ 小的原子; $W_n$ 是由 $\psi$ 定义的加性鞅 (见附录); $\xi$ 表示由\nameref{lemma:many-to-one}给出分布的随机变量, 称为脊柱随机游走步长, 其方差为 $\sigma^2:=\psi''(1)$; $\ell_N$ 是以 $N$ 为索引的时间序列; $v_N$ 为最左侧粒子的渐近速率. 假设的详细说明见下一章第 \ref{section:explanations-on-assumptions} 节:
    }
    \begin{assumption}[\bilingual{Supercritical Regime}{超临界情形}]
        \label{assumption:psi-zero}
        \bilingual{
            $\psi(0)>0$. That is,~ $\mathbb E[\#\Xi]>1$.
        }{
            $\psi(0)>0$. 即 $\mathbb E[\#\Xi]>1$.
        }
    \end{assumption}
    \begin{assumption}[\bilingual{At Least One Child}{至少一个子代}]
        \label{assumption:at-least-one-child}
        \bilingual{
            $\#\Xi\ge 1$ a.s. Or equivalently,~ $\mathbb P(\#\Xi=0)=0$.
        }{
            $\#\Xi\ge 1$ a.s. 或者等价地,~ $\mathbb P(\#\Xi=0)=0$.
        }
    \end{assumption}
    \begin{assumption}[\bilingual{Exponential Moment}{负指数矩}]
        \label{assumption:psi-one}
        \bilingual{
            $\psi(1)=0$. That is,~ $\mathbb E[W_1]=\mathbb E\left[\sum_{x\in\Xi}e^{-x}\right]=1$.
        }{
            $\psi(1)=0$. 即 $\mathbb E[W_1]=\mathbb E\left[\sum_{x\in\Xi}e^{-x}\right]=1$.
        }
    \end{assumption}
    \begin{assumption}[\bilingual{Zero Mean}{零均值}]
        \label{assumption:psi-prime-one}
        \bilingual{
            $\psi'(1)=0$. That is,~ for the spine random walk step $\xi$,~ the mean $\mathbb E[\xi]=0$.
        }{
            $\psi'(1)=0$. 即,~ 对于脊柱随机游走步长$\xi$,~ 均值$\mathbb E[\xi]=0$.
        }
    \end{assumption}
    \begin{assumption}[\bilingual{Finite Variance}{有限方差}]
        \label{assumption:finite-variance}
        \bilingual{
            $\psi''(1)<\infty$. That is,~ for the step $\xi$,~ the variance $\mathbb E[\xi^2]=\sigma^2<\infty$.
        }{
            $\psi''(1)<\infty$. 即,~ 对于步长$\xi$,~ 方差$\mathbb E[\xi^2]=\sigma^2<\infty$.
        }
    \end{assumption}
    \begin{assumption}[\bilingual{Finite Cross Log-Laplacian}{有限交叉Laplace变换}]
        \label{assumption:log-expt-cross-sum-exp-neg-sum}
        {
            \setlength\belowdisplayskip{0pt}
            \begin{equation*}
                \log\mathbb E\!\left[
                    \sum_{
                        i,j\in\mathbb N,\,
                        i\neq j
                    }
                    e^{
                        -(\Xi_i+\Xi_j)
                    }
                \right]
                <
                \infty.
            \end{equation*}
        }
    \end{assumption}
    \begin{assumption}[\bilingual{First Moment of Positive Part of Leftmost Step}{最左侧正部一阶矩}]
        \label{assumption:moment-1}
        \bilingual{
            Under \Cref{assumption:at-least-one-child}, we further assume that
        }{
            在\Cref{assumption:at-least-one-child} 下,~ 进一步假设
        }
        $\mathbb E[(\Xi_1)_+]<\infty$.
    \end{assumption}
    \begin{assumption}[\bilingual{Fourth Moment of Positive Part of Leftmost Step}{最左侧正部四阶矩}]
        \label{assumption:moment-4}
        \bilingual{
            Under \Cref{assumption:moment-1}, we further assume that
        }{
            在\Cref{assumption:moment-1} 下,~ 进一步假设
        }
        $\mathbb E[(\Xi_1)_+^4]<\infty$.
    \end{assumption}

    \bilingual{
        With the above assumptions, the main conclusions obtained in this paper are as follows:
    }{
        本文得到的主要结论如下:
    }

    \begin{theorem}
        \label{proposition:3.1}
        \bilingual{
            Fix $a\in(0,1]$ and $t>0$. If $\mathcal V^N=\mathcal V^{N|a}$ starts with $\lfloor N^a\rfloor$ particles located at $0$, and its branching law $\Xi$ satisfies \Cref{assumption:psi-zero,assumption:at-least-one-child,assumption:psi-one,assumption:psi-prime-one,assumption:finite-variance,assumption:log-expt-cross-sum-exp-neg-sum,assumption:moment-1,assumption:moment-4}, then for any sequence $\ell_N$ satisfying $\lim_{N\to\infty}\frac{\ell_N}{\log^3N}=t$,
        }{
            固定 $a\in(0,1]$ 和 $t>0$. ~若 $N$-分支随机游走 $\mathcal V^N=\mathcal V^{N|a}$ 从 $\lfloor N^a\rfloor$ 个位于 $0$ 的粒子开始,~ 且其分支律$\Xi$满足\Cref{assumption:psi-zero,assumption:at-least-one-child,assumption:psi-one,assumption:psi-prime-one,assumption:finite-variance,assumption:log-expt-cross-sum-exp-neg-sum,assumption:moment-1,assumption:moment-4},~ 那么对于任何满足 $\displaystyle\lim_{N\to\infty}\textstyle\frac{\ell_N}{\log^3N}=t$ 的数列 $\ell_N$,
        }
        {
            \begin{equation*}
                \frac 1{\log N}
                \min_{|u|=\ell_N}
                \mathcal V^N(u)
                \xrightarrow{\mathrm L^2}
                -a+\frac{\pi^2}2\sigma^2t.
            \end{equation*}
        }
        \bilingual{
            Moreover, under an additional \Cref{assumption:moment-exp-1}, i.e., the existence of $\lambda>0$ such that $\mathbb E[e^{\lambda \Xi_1}]<\infty$, the $\mathrm L^2$ convergence holds for $a=0$.
        }{
            另外,~ 在额外的\Cref{assumption:moment-exp-1} 下,~ 即存在 $\lambda>0$ 使得 $\mathbb E[e^{\lambda \Xi_1}]<\infty$,~ 该$\mathrm L^2$ 收敛对于 $a=0$ 也成立.
        }
    \end{theorem}

    \begin{theorem}
        \label{theorem:speed}
        \bilingual{
            For $a\in[0,1]$, consider a branching random walk $\mathcal V^{N|a}$ that starts with $\lfloor N^a\rfloor$ particles at position $0$, and whose branching law $\Xi$ satisfies \Cref{assumption:psi-zero,assumption:at-least-one-child,assumption:psi-one,assumption:psi-prime-one,assumption:finite-variance,assumption:log-expt-cross-sum-exp-neg-sum,assumption:moment-1,assumption:moment-4}. Then the speed of the leftmost particle in $\mathcal V^{N|a}$ exists and is given by
        }{
            对于 $a\in[0,1]$,~ 考虑从位置 $0$ 处的 $\lfloor N^a\rfloor$ 个粒子出发的分支随机游走 $\mathcal V^{N|a}$,~ 其分支律 $\Xi$ 满足\Cref{assumption:psi-zero,assumption:at-least-one-child,assumption:psi-one,assumption:psi-prime-one,assumption:finite-variance,assumption:log-expt-cross-sum-exp-neg-sum,assumption:moment-1,assumption:moment-4}. ~则 $\mathcal V^{N|a}$ 中最左侧粒子的速率存在且由下式给出:
        }
        {
            \setlength\belowdisplayskip{-10pt}
            \begin{equation*}
                v_N
                =
                \lim_{n\to\infty}
                \frac{1}{n}\mathcal M^{N|a}_n
                =
                \frac{\pi^2}2\sigma^2\frac{1}{\log^2 N}
                +
                o_N\left(\frac{1}{\log^2 N}\right).
            \end{equation*}
        }
    \end{theorem}

    \begin{theorem}[\bilingual{Speed under a First-Moment Assumption}{一阶矩假设下的渐近速率}]
        \label{theorem:speed-L1}
        \bilingual{
            Consider the $N$-branching random walk $\mathcal V^{N|a}$ starting with $\lfloor N^a\rfloor$ particles at position $0$, where $a\in[0,1]$. Suppose that its branching law $\Xi$ satisfies \Cref{assumption:psi-zero,assumption:at-least-one-child,assumption:psi-one,assumption:psi-prime-one,assumption:finite-variance,assumption:log-expt-cross-sum-exp-neg-sum,assumption:moment-1}. Then the almost-sure speed of its leftmost particle exists and
        }{
            考虑从位置 $0$ 处的 $\lfloor N^a\rfloor$ 个粒子出发的 $N$-分支随机游走 $\mathcal V^{N|a}$,~ 其中 $a\in[0,1]$. ~设其分支律 $\Xi$ 满足\Cref{assumption:psi-zero,assumption:at-least-one-child,assumption:psi-one,assumption:psi-prime-one,assumption:finite-variance,assumption:log-expt-cross-sum-exp-neg-sum,assumption:moment-1}. ~则最左侧粒子的几乎处处渐近速率存在,~ 且
        }
        \begin{equation*}
            v_N
            =\lim_{n\to\infty}\frac{\mathcal M^{N|a}_n}{n}
            =\frac{\pi^2\sigma^2}{2(\log N)^2}
            +o_N\!\left((\log N)^{-2}\right).
        \end{equation*}
    \end{theorem}

    \bilingual{
        \emph{Proof status of \Cref{theorem:speed-L1}.} The one-sided $L^1$ estimates needed to remove \Cref{assumption:moment-4} have not been established in this chapter. This stronger statement is recorded as a proposed extension; the speed theorem proved here is \Cref{theorem:speed}.
    }{
        \emph{\Cref{theorem:speed-L1} 的证明状态.} ~本章尚未建立去掉\Cref{assumption:moment-4} 所需的单侧 $L^1$ 估计. ~此处将该加强结论作为待证的推广版本记录;~ 本文已证明的速率定理为\Cref{theorem:speed}.
    }

    \bilingual{
        The rest of this paper is organized as follows: Chapter 2 will formally define $N$-branching random walks in Section \ref{section:formalization-and-branching-properties}, and provide detailed explanations of the above assumptions in Section \ref{section:explanations-on-assumptions}, with some lemmas on the assumptions proved in Appendix \ref{section:lemmas-on-assumptions}. The other parts of Chapter 2 will explain the research tools and some known results used in this paper, including \nameref{lemma:branching-property}, \nameref{lemma:many-to-one} and \nameref{theorem:mogulskii}. The four sections of Chapter 3 establish preparatory estimates, where Section \ref{section:mogulskii-lemmas} is the lemma proof applying \nameref{theorem:mogulskii}, Section \ref{section:spine-random-walk-lemmas} further applies \nameref{theorem:mogulskii} to the spine random walk to obtain its properties, Section \ref{section:lemmas-on-brw} further uses \nameref{lemma:many-to-one} to obtain lemmas on branching random walks, and Section \ref{section:lemmas-on-n-brw} gives the coupling properties of $N$-branching random walks and argues for the existence of asymptotic speed. Chapter 4 proves the trajectory limit in Section \ref{section:main-results} and treats the speed in Section \ref{section:speed-proof}. Chapter 5 is conclusions and future work.
    }{
        第二章第 \ref{section:formalization-and-branching-properties} 节将形式化定义 $N$-分支随机游走,~ 第 \ref{section:explanations-on-assumptions} 节将对上述假设进行详细解释说明, 部分关于假设的引理证明在附录 \ref{section:lemmas-on-assumptions},~ 第二章的其他部分则说明了本文所用的研究工具和一些已知结论,~ 包括\nameref{lemma:branching-property}、\nameref{lemma:many-to-one} 和\nameref{theorem:mogulskii}. 第三章的四节建立预备估计, 其中第 \ref{section:mogulskii-lemmas} 节是应用\nameref{theorem:mogulskii}的引理证明,~ 第 \ref{section:spine-random-walk-lemmas} 对脊柱分解得到的随机游走进一步应用\nameref{theorem:mogulskii}得到它们的性质,~ 第 \ref{section:lemmas-on-brw} 节进一步利用\nameref{lemma:many-to-one}得到分支随机游走的引理,~ 第 \ref{section:lemmas-on-n-brw} 节则给出$N$-分支随机游走的耦合性质,~ 并论证了渐近速度的存在性. ~第四章第 \ref{section:main-results} 节讨论轨迹极限,~ 第 \ref{section:speed-proof} 节讨论渐近速率. ~第五章是结论与展望.
    }
